\begin{table}[H]
  \centering
  \footnotesize
  \begin{tabularx}{\textwidth}{@{}
      >{\raggedright\arraybackslash}p{0.24\textwidth}
      >{\raggedright\arraybackslash}X
      >{\raggedright\arraybackslash}p{0.12\textwidth}@{}}
    \toprule
    Name & One-line claim & Substrate cite \\
    \midrule
    Quotientality theorem (\Fref{F10}) &
    Every observable constant on access classes factors uniquely through the access quotient, and the access quotient factors through every surjection on whose fibers the signature is constant (it is the coarsest such quotient).  &
    \textemdash \\
    Adequacy / No-Overread theorem (\Fref{F11}) &
    Under the no-overread rule, accepted exact claims factor uniquely through the access quotient.  &
    \textemdash \\
    No-Free-Distinction theorem (\Fref{F12}) &
    A readout is current-visible iff it factors through the access
    quotient; the pairing with it is always the coarsest map through which
    both factor, strict when the readout splits a fiber, and then, under
    (NO), the readout is not accepted as exact.  &
    \textemdash \\
    Hiddenness Normal Form (\Fref{F13a}), conditional &
    Under the source hypothesis \((\mathsf S)\) (surjectivity and four of
    the five conjuncts), the fifth conjunct holds, and exposure, collapse
    and absence of surplus are equivalent; without \((\mathsf S)\),
    predictive surplus holds exactly when predictive collapse fails, for a
    surjective current quotient. &
    \PaperPvNP \\
    CSL Formed-Closure (\Fref{F13b}) &
    With a determining-state readout, every predictive-surplus pair is separated by the determining state.  &
    \PaperCSL \\
    Layer-Multiplicity Theorem (\Fref{F24}) &
    A role readout splits an access quotient exactly when its role
    obstruction is nonempty, and a role split implies the strict
    layer extension by the joint role-layer quotient.  &
    \textemdash \\
    \bottomrule
  \end{tabularx}
  \caption{Access at a glance: the six laws of \Cref{sec:access}.}
  \label{tab:axis-access-summary}
\end{table}
